\documentclass[11pt]{article}
\usepackage[utf8]{inputenc}
\usepackage[T1]{fontenc}
\usepackage{lmodern}
\usepackage[margin=2.4cm]{geometry}
\usepackage{amsmath,amssymb}
\usepackage{booktabs}
\usepackage{array}
\usepackage{enumitem}
\usepackage{caption}
\usepackage{xcolor}
\usepackage{titlesec}
\usepackage[hidelinks]{hyperref}

\definecolor{ink}{HTML}{23211D}
\definecolor{accent}{HTML}{8C7B66}
\definecolor{rulelt}{HTML}{D9D2C7}

\titleformat{\section}{\sffamily\large\bfseries\color{ink}}{\thesection.}{0.6em}{}[{\color{rulelt}\titlerule[0.6pt]}]
\titleformat{\subsection}{\sffamily\normalsize\bfseries\color{ink}}{\thesubsection}{0.6em}{}
\titlespacing*{\section}{0pt}{2.1ex plus 1ex minus .2ex}{1.0ex}
\captionsetup{font=small,labelfont=bf,labelsep=period}
\setlength{\parindent}{0pt}
\setlength{\parskip}{0.45em}

\newcommand{\E}{\mathbb{E}}
\newcommand{\R}{\mathbb{R}}
\newcommand{\T}{^{\!\top}}
\newenvironment{where}{\par\nobreak\vspace{-0.2em}\begin{itemize}[leftmargin=1.5em,itemsep=1pt,topsep=2pt,parsep=0pt,label=\textcolor{accent}{\textbullet}]\small}{\end{itemize}\vspace{0.2em}}

\title{\vspace{-1.2cm}\sffamily\bfseries\color{ink} Methodological Note (v5)\\[2pt]
\large Replication of Boucher, Rendall, Ushchev and Zenou (2024)\\
\large\itshape Toward a General Theory of Peer Effects}
\author{Juan Valderrama}
\date{June 2026}

\begin{document}
\maketitle
\vspace{-0.7cm}

\noindent\textbf{Abstract.} This note documents the consolidated final code (\emph{v5}) for the replication. It explains, with the math worked out and every symbol defined, the structural model, why estimation uses GMM, where identification comes from, the concentrated two-step estimator and the continuously updated estimator (CUE), and the full battery of diagnostics. It adds two things relative to earlier drafts: a comparison of several instrument transformations (to see which identify $\beta$ worse) and the over-identification J-test. All figures in the results come from the runs with 150 schools and 30{,}000 students. The estimation is clean because the data come from a correctly specified synthetic DGP; this validates method and code, not an empirical claim about peer effects.

% =====================================================================
\section{Goal and strategy}
The paper generalizes the \emph{linear-in-means} (LIM) peer model, where each agent is influenced by the \emph{average} of their peers, to a CES social norm with a curvature parameter $\beta$ that decides \emph{which} peers one responds to. With no access to the real microdata (Add Health), we run a Monte Carlo: (i) generate synthetic data \emph{from the model} with known true parameters; (ii) estimate; (iii) check recovery. If we know the truth, we can validate the estimator and the code.

\subsection{Notation}
\begin{table}[h]\centering\small
\caption{Symbols used throughout.}
\begin{tabular}{ll}
\toprule
\multicolumn{2}{l}{\textbf{Model}}\\
\midrule
$i,\,j$ & student indices ($j$: friends of $i$); $s(i)$ school of $i$ \\
$y_i$ & observed GPA (the outcome) of student $i$ \\
$p_i$ & private component: $i$'s grade absent peer influence \\
$x_i,\,\gamma,\,c$ & covariates (age, female, college parent), their coefficients, intercept \\
$u_{s(i)},\,\varepsilon_i$ & school fixed effect; idiosyncratic error, $\E[\varepsilon_i\mid x,G,u]=0$ \\
$g_{ij},\,G$ & weight $i$ gives friend $j$ (row-normalized); network matrix \\
$S_i(\beta)$ & CES social norm of $i$'s peers \\
$\beta$ & curvature of the norm (which peers count) \\
$\lambda,\,\delta$ & peer-effect intensity; weight on the private component \\
$\lambda_1=\lambda+\delta-1$ & endogenous spillover (contagion) \\
$\lambda_2=1-\delta$ & conformism \\
\midrule
\multicolumn{2}{l}{\textbf{Estimation}}\\
\midrule
$\widehat y_i$ & first-stage predicted GPA (own covariates + school FE) \\
$\widehat S_i$ & norm at $\widehat y$: instrument for the level of the norm \\
$\widehat D_i=\partial\widehat S_i/\partial\beta$ & instrument that identifies $\beta$ \\
$Z_i,\,r_i(\theta)$ & instrument vector; structural residual \\
$\theta=(\gamma,\lambda,\beta,\delta)$ & parameters; $\bar g(\theta)=\tfrac1n\sum_i Z_i r_i(\theta)$ moments \\
$W,\,\Omega$ & GMM weight matrix; variance of the moments \\
$Q(\theta)$ & GMM objective $\bar g\T W \bar g$ \\
$J(\beta_0),\,J=nQ$ & Anderson--Rubin statistic; Hansen over-identification statistic \\
$n,\,G_{\text{cl}}$ & students $(30{,}000)$; schools / clusters $(150)$ \\
\bottomrule
\end{tabular}
\end{table}

% =====================================================================
\section{The structural model}
For each student $i$,
\begin{align}
\text{isolated:}\quad & y_i = p_i, &
\text{connected:}\quad & y_i = \delta\,p_i + \lambda\,S_i(\beta),
\end{align}
\begin{where}
\item $y_i$ --- GPA; $p_i$ --- private component; $S_i(\beta)$ --- peer norm.
\item $\delta$ --- weight on own private component; $\lambda$ --- weight on peers.
\item isolated $=$ no friends (no peer term); connected $=$ at least one friend.
\end{where}
The private component is
\begin{equation}
p_i = c + x_i\T\gamma + u_{s(i)} + \varepsilon_i, \qquad \E[\varepsilon_i\mid x,G,u]=0 .
\end{equation}
\begin{where}
\item $c$ --- common intercept; $x_i\T\gamma$ --- covariate part; $u_{s(i)}$ --- school fixed effect.
\item $\varepsilon_i$ --- mean-zero idiosyncratic error given observables.
\end{where}

\subsection{What $\beta$ does: a power mean}
\begin{equation}
S_i(\beta)=\Bigl(\sum_j g_{ij}\,y_j^{\,\beta}\Bigr)^{1/\beta},\qquad \sum_j g_{ij}=1 .
\end{equation}
\begin{where}
\item $g_{ij}$ --- friendship weights (sum to 1 $\Rightarrow$ weighted average); $y_j$ --- friends' grades.
\item $\beta$ --- curvature: $\beta=1$ mean (LIM); $\beta>1$ over-weights high achievers; $\beta\to\infty$ the max.
\end{where}
LIM is the special case $\beta=1$; the empirical question is which $\beta$ fits (here the true value is $10$).

\subsection{Reparametrization and equilibrium}
\begin{equation}
\lambda_1=\lambda+\delta-1,\quad \lambda_2=1-\delta,\qquad
y_i = p_i + \underbrace{\lambda_1 S_i}_{\text{contagion}} + \underbrace{\lambda_2(S_i-p_i)}_{\text{conformism}} .
\end{equation}
Because $y_i$ depends on peers' $y_j$, the school's outcomes solve a fixed point $y=T(y)$ with $T(y)=\delta p+\lambda S(y;\beta)$. The CES norm is homogeneous of degree~1, its Jacobian has row sums $\le 1$, so $T$ is a contraction with modulus $\le\lambda<1$: a unique equilibrium, computed by iteration ($\approx 16$ steps). At $\beta=1$, $y=(I-\lambda G)^{-1}\delta p$ and a uniform shock is amplified by the \emph{social multiplier} $\delta/(1-\lambda)=0.85/0.75=1.133$.
\begin{where}
\item $\lambda_1$ --- pure spillover; $\lambda_2$ --- pressure toward the group norm; $\lambda_1>0\iff\delta+\lambda>1$.
\item $(I-\lambda G)^{-1}=I+\lambda G+\lambda^2G^2+\cdots$ --- sum of all feedback rounds.
\end{where}
% =====================================================================
\section{The econometric problem}
\subsection{Reflection and endogeneity}
$S_i$ contains the friends' $y_j$, which depend on $\varepsilon_j$ and, through the equilibrium feedback, on $\varepsilon_i$. Hence
\begin{equation}
\E[S_i\,\varepsilon_i]\neq 0 ,
\end{equation}
so OLS of $y_i$ on $S_i$ is inconsistent: $\hat\lambda_{\text{OLS}}\xrightarrow{p}\lambda+\operatorname{Cov}(S_i,\varepsilon_i)/\operatorname{Var}(S_i)$. This is Manski's (1993) reflection problem plus simultaneity.

\subsection{Why GMM (not OLS or MLE)}
GMM (Hansen, 1982) needs only \emph{instruments} $Z_i$ with $\E[Z_i\,r_i(\theta_0)]=0$. We use it because $\beta$ enters \emph{nonlinearly} (OLS cannot) and because GMM is semiparametric (no distributional assumption on $\varepsilon$, unlike MLE) and robust to clustering; with optimal weights it is efficient.

% =====================================================================
\section{Identification and instruments}
\subsection{Network intransitivities and isolated students}
Bramoull\'e, Djebbari and Fortin (2009): in a triad $A\!\to\!B\!\to\!C$ with $A\!\not\to\!C$, $C$'s characteristics shift $A$'s norm (through $B$) but do not enter $A$'s equation, so they instrument $S_A$. Network gaps (intransitivities) create the instruments. Separately, isolated students obey $y_i=c+x_i\T\gamma+u+\varepsilon_i$ with no endogeneity, so they pin down $\gamma$ directly (hence $\gamma$ is very precise).

\subsection{The baseline instruments $\widehat S$ and $\widehat D$}
The first stage predicts $\widehat y$ from own covariates and school effects, then builds
\begin{equation}
\widehat S_i=\Bigl(\sum_j g_{ij}\widehat y_j^{\,\beta}\Bigr)^{1/\beta},\qquad
\widehat D_i=\frac{\partial\widehat S_i}{\partial\beta}
=\widehat S_i\Bigl[\tfrac1\beta\tfrac{\sum_j g_{ij}\widehat y_j^{\beta}\ln\widehat y_j}{\sum_j g_{ij}\widehat y_j^{\beta}}-\tfrac1{\beta^2}\ln\!\textstyle\sum_j g_{ij}\widehat y_j^{\beta}\Bigr].
\end{equation}
\begin{where}
\item $\widehat y_j$ --- friend's predicted grade (function of exogenous $x$ only).
\item $\widehat S_i$ --- instruments the \emph{level} of the norm $\Rightarrow$ identifies $\lambda$.
\item $\widehat D_i=\partial\widehat S_i/\partial\beta$ --- instruments the \emph{curvature} $\Rightarrow$ identifies $\beta$. It is tiny at high $\beta$ (the norm saturates), which is why $\beta$ is the hard parameter.
\item The baseline does \emph{not} use friend-of-friend covariates ($G^2X$): the first stage uses only own covariates, so $\widehat S,\widehat D$ already encode friends' (distance-1) characteristics. $G^2X$ is offered as an option (below).
\end{where}
The non-isolated moments stack $Z_i=(x_i,\widehat S_i,\widehat D_i)$ (5 moments) and the isolated $Z_i=x_i$ (3): 8 moments for 6 parameters, over-identified by~2.

\subsection{Configurable instruments [v5]}
v5 makes the non-isolated instrument set a switch (\texttt{ACTIVE\_INSTRUMENTS}). Beyond the baseline, one can append $\widehat D^2$, $\widehat S^2$, the norm at $\beta=1$, friends' covariates $GX$, or friends-of-friends $G^2X$, and \texttt{compare\_instrument\_sets} refits under each to show which estimate $\beta$ worse (Section~\ref{sec:instr}).

% =====================================================================
\section{The estimator}
\subsection{Concentrated two-step GMM}
\textbf{(a)} Within transform by (school $\times$ isolation) removes $u_{s(i)}$ and $c$.
\textbf{(b)} Given $(\lambda,\beta,\delta)$, $\gamma$ is linear, solved analytically, so the search is over 3 parameters:
\begin{equation}
Q(\lambda,\beta,\delta)=\bar g(\theta)\T W\,\bar g(\theta),\qquad \hat\theta=\arg\min_\theta Q(\theta).
\end{equation}
\textbf{(c)} Two steps: $W=I$ first ($\hat\theta_1$), then $\widehat W=\widehat\Omega^{-1}$ with $\widehat\Omega=\tfrac1n\sum_i Z_iZ_i\T r_i(\hat\theta_1)^2$, re-optimized for the efficient $\hat\theta_2$. A global $\beta$ prescan plus multistart avoid spurious optima.
\begin{where}
\item $W$ --- weight matrix; $\widehat W=\widehat\Omega^{-1}$ is the efficient choice.
\end{where}

\subsection{CUE [v3 / v5]}
The continuously updated estimator recomputes the weight matrix at \emph{every} $\theta$ inside the objective:
\begin{equation}
Q_{\text{CUE}}(\theta)=\bar g(\theta)\T\,\Omega(\theta)^{-1}\,\bar g(\theta).
\end{equation}
\begin{where}
\item $\Omega(\theta)$ --- moment variance recomputed at each $\theta$ (vs fixed in two-step).
\item Motivation: with a weak instrument CUE can in principle have better finite-sample behavior. Finding (Section~\ref{sec:cue}): in this correctly specified DGP it yields \emph{no practical improvement} over the two-step estimator.
\end{where}

% =====================================================================
\section{Inference}
\textbf{Cluster-robust SE} (sandwich, school clusters, with $G_{\text{cl}}/(G_{\text{cl}}-1)$):
\begin{equation}
\widehat V=(\widehat G\T W\widehat G)^{-1}\widehat G\T W\,\widehat\Lambda\,W\widehat G(\widehat G\T W\widehat G)^{-1},\quad
\widehat\Lambda=\tfrac{G_{\text{cl}}}{G_{\text{cl}}-1}\sum_s\Bigl(\textstyle\sum_{i\in s}Z_ir_i\Bigr)\Bigl(\textstyle\sum_{i\in s}Z_ir_i\Bigr)\T .
\end{equation}
\begin{where}
\item $\widehat G=\partial\bar g/\partial\theta\T$ --- moment Jacobian; $\widehat\Lambda$ --- school-clustered moment variance.
\item $\lambda_1,\lambda_2$ SEs by the delta method; clustering reflects $150$ schools, not $30{,}000$ independent students.
\end{where}
\textbf{Anderson--Rubin} (weak-IV robust): fix $\beta=\beta_0$, profile the rest, $J(\beta_0)=n\,Q(\beta_0)\xrightarrow{d}\chi^2_3$; the 95\% set is $\{\beta_0:J(\beta_0)\le 7.81\}$.
\textbf{Hansen $J$-test [v5]} (the extra moments as a test): at the efficient step $J=n\,Q(\hat\theta)\xrightarrow{d}\chi^2_{\,8-6}=\chi^2_2$; a large $J$ rejects the over-identifying restrictions.
\textbf{LIM test:} compare the objective at $\beta=1$ to the free optimum.
% =====================================================================
\section{Results}
\subsection{Generated data}
\begin{table}[h]\centering\small
\caption{Synthetic environment (DGP\_v5).}
\begin{tabular}{lr@{\hspace{2.2em}}lr}
\toprule
Schools & 150 & Mean GPA & 2.778 \\
Students & 30{,}000 & Min / Max GPA & 1.378 / 4.000 \\
Isolated & 6{,}080 (20.3\%) & Clipped to $[1,4]$ & 1 (0.003\%) \\
\bottomrule
\end{tabular}
\end{table}

\subsection{Parameter recovery}
\begin{table}[h]\centering\small
\caption{Two-step and CUE estimates (150 schools). School-cluster robust SE; $t$ for the two-step.}
\begin{tabular}{lrrrrr}
\toprule
Parameter & True & Two-step & CUE & Cluster SE & $t$ \\
\midrule
$\gamma_{\text{age}}$ & 0.10 & 0.0985 & 0.0985 & 0.0014 & 71.4 \\
$\gamma_{\text{female}}$ & 0.24 & 0.2401 & 0.2401 & 0.0032 & 75.2 \\
$\gamma_{\text{parent coll.}}$ & 0.40 & 0.3984 & 0.3984 & 0.0045 & 89.2 \\
$\lambda$ & 0.25 & 0.2518 & 0.2517 & 0.0076 & 33.2 \\
$\beta$ & 10.00 & 9.350 & 9.359 & 1.105 & 8.46 \\
$\delta$ & 0.85 & 0.8486 & 0.8486 & 0.0091 & 92.9 \\
$\lambda_1$ & 0.10 & 0.1004 & 0.1003 & 0.0116 & 8.66 \\
$\lambda_2$ & 0.15 & 0.1514 & 0.1514 & 0.0091 & 16.57 \\
\midrule
GMM objective $Q$ & --- & 0.000185 & 0.000185 & & \\
\bottomrule
\end{tabular}
\end{table}
$\gamma,\lambda,\delta$ land almost exactly on the truth; $\beta$ is the least precise. Two-step and CUE coincide at this design.

\subsection{First stage, instruments and tests}
\begin{table}[h]\centering\small
\caption{Instrument strength and weak-IV / over-identification / LIM tests.}
\begin{tabular}{lr@{\hspace{2.0em}}lr}
\toprule
Partial $F$ (for $\beta$) & 38.5 & AR 95\% set for $\beta$ & $[9.0,\,10.5]$ \\
First-stage $R^2$ ($\widehat y$) & 0.575 & $\chi^2_{0.95,3}$ & 7.81 \\
$\mathrm{corr}(S,\widehat S)$ & 0.701 & Hansen $J=nQ$ (dof 2) & 5.56 ($p=0.062$) \\
mean $\widehat D$ (levels) & 0.0041 & LIM obj.\ vs free & 0.0071 vs 0.00019 \\
\bottomrule
\end{tabular}
\end{table}
The AR set contains the true $\beta=10$. The $J$-test does not reject the model (close, $p=0.06$, single draw). The LIM ($\beta=1$) is rejected: the distance statistic $n(Q_{\text{LIM}}-Q_{\text{free}})\approx 209$ on $\chi^2_1$ gives $p\approx 0$ --- the norm is not the simple mean.

\subsection{Instrument transformations}\label{sec:instr}
\begin{table}[h]\centering\small
\caption{Refitting under different instrument sets (illustrative, 30 schools, true $\beta=10$). Read: $\beta$ far from 10, huge $\mathrm{SE}$, or tiny $F$ $\Rightarrow$ worse.}
\begin{tabular}{lrrr}
\toprule
Instrument set & $\widehat\beta$ & $\mathrm{SE}(\beta)$ & $F$ \\
\midrule
baseline $[X,\widehat S,\widehat D]$ & 8.85 & 2.22 & 8.9 \\
$+\,\widehat D^2$ (the trap) & 8.10 & 2.08 & 8.7 \\
$+\,\widehat S^2$ & 10.73 & 3.10 & 8.9 \\
$+\,\widehat S(\beta{=}1)$ & 8.39 & 2.04 & 8.8 \\
$+\,GX$ (friends) & 8.44 & 2.04 & 8.8 \\
$+\,G^2X$ (friends of friends) & 9.38 & 2.15 & 9.0 \\
$\widehat S$ only (no $\widehat D$) & 2.00 & 57.2 & 0.5 \\
\bottomrule
\end{tabular}
\end{table}
The qualitative ranking is robust: dropping $\widehat D$ leaves $\beta$ \emph{unidentified} ($\mathrm{SE}=57$), confirming $\widehat D$ carries the curvature signal; extra derivative-like moments ($\widehat D^2$, $\widehat S^2$) do not help and worsen the fit; friend / friend-of-friend covariates are comparable to the baseline.

\subsection{Two-step vs CUE: no practical improvement}\label{sec:cue}
\begin{table}[h]\centering\small
\caption{Monte Carlo, two-step vs CUE, true $\beta=10$, two designs.}
\begin{tabular}{llrrrr}
\toprule
Design & Estimator & mean $\widehat\beta$ & median & RMSE & coverage \\
\midrule
30 schools (21 reps) & Two-step & 10.50 & 9.73 & 2.76 & 0.95 \\
                     & CUE      & 10.55 & 9.73 & 2.78 & 0.95 \\
\midrule
150 schools (8 reps) & Two-step & \phantom{0}9.91 & 9.82 & 0.54 & 1.00 \\
                     & CUE      & \phantom{0}9.92 & 9.84 & 0.53 & 1.00 \\
\bottomrule
\end{tabular}
\end{table}
The two estimators are statistically indistinguishable at both designs. At 150 schools (the real design) they differ by less than $0.05$ in $100\%$ of draws (mean $|\widehat\beta_{\text{CUE}}-\widehat\beta_{\text{2s}}|=0.013$), with identical RMSE, SE and coverage; at 30 schools their medians, RMSE and tails coincide. An earlier 50-draw run had appeared to favour CUE on RMSE, but that gap was driven by a single two-step blow-up in one seed set and does not survive re-sampling. \textbf{Conclusion:} CUE recomputes $\Omega(\theta)^{-1}$ at every $\theta$ at extra computational cost and delivers \emph{no} improvement here; we keep it as a documented option and report the two-step estimator (with the Anderson--Rubin set for $\beta$) as primary.

\subsection{Monte Carlo and identification frontier}
\begin{table}[h]\centering\small
\caption{Left: bias / 95\% coverage (36 reps, 30 schools). Right: $\beta$ frontier.}
\begin{tabular}{lrr@{\hspace{3em}}rrr}
\toprule
Param. & Bias & Cov. & Schools & $\mathrm{SE}(\beta)$ & $F$ \\
\midrule
$\lambda$ & $-0.004$ & 0.94 & 20 & 3.14 & 5.9 \\
$\beta$ & $+0.351$ & 0.97 & 40 & 2.70 & 12.0 \\
$\delta$ & $+0.001$ & 0.97 & 80 & 1.78 & 23.0 \\
$\lambda_1$ & $-0.003$ & 0.94 & 150 & 1.09 & 37.1 \\
$\lambda_2$ & $-0.001$ & 0.97 & 300 & 0.85 & 78.9 \\
\bottomrule
\end{tabular}
\end{table}
$\mathrm{SE}(\beta)$ falls and $F$ rises with the number of schools: more (and denser) network data is the lever for $\beta$.

% =====================================================================
\section{Why the estimation comes out clean}
Plainly, the experiment is built to be clean, not because peer effects are easy to measure:
\begin{enumerate}\itemsep2pt
\item \textbf{Zero misspecification:} data come from exactly the estimated model ($Q\approx0$).
\item \textbf{Instruments valid by construction:} network and errors independent; no homophily.
\item \textbf{No measurement error}, large $n$, the only confounder (school) swept exactly, and isolated students pin down $\gamma$.
\end{enumerate}
The only genuine difficulty is $\beta$: at high $\beta$ the CES norm saturates and $\widehat D\to 0$. This validates \textbf{method and code}; with real data it would be messier.

% =====================================================================
\section{Conclusion}
The pipeline is: a CES peer model with a fixed-point equilibrium; endogeneity (reflection) that rules out OLS; identification from network intransitivities plus the isolated subsample; concentrated two-step GMM and CUE with $\gamma$ profiled out and cluster-robust SE; and a battery of diagnostics (recovery, Anderson--Rubin, Hansen $J$, LIM, instrument-transformation comparison, two-step-vs-CUE Monte Carlo, identification frontier). All eight parameters are recovered with near-nominal coverage; $\beta$ is consistent but weakly identified at high values, a fact the diagnostics quantify. The consolidated code lives in the \texttt{*\_v5} modules; \texttt{python run\_all\_v5.py} reproduces everything in one pass.

\end{document}
